\subsection{惯性定律。}

定理 1（“惯性定律”）. --- 假设 A 满足本节开头的诸假设，且 E 的维数为有限数 n。设 Φ 是 E 上的一个埃尔米特型。则：

\begin{enumerate}
\def\labelenumi{\alph{enumi})}
\item
  存在 E 的一个分解，把它写成与 E 正交的子空间 \(E^{0}\) 与两个子空间 \(E^{+}\) 与 \(E^{-}\) 的直和，使得 \(\Phi\) 在 \(E^{+}\) （相应地 \(E^{-}\) ）上的限制是正的（相应地，负的）且非退化。
\item
  存在 E 的正交基 \((e_i)_{1\leqslant i\leqslant n}\) ，使得
\end{enumerate}

\[
\Phi (\sum_ {i = 1} ^ {n} \xi_ {i} e _ {i}, \sum_ {i = 1} ^ {n} \eta_ {i} e _ {i}) = \sum_ {i = 1} ^ {s} \xi_ {i} \overline {{\eta}} _ {i} - \sum_ {i = s + 1} ^ {s + t} \xi_ {i} \overline {{\eta}} _ {i}.\tag{2}
\]

\begin{enumerate}
\def\labelenumi{\alph{enumi})}
\setcounter{enumi}{2}
\item
  \(E^{+}\) 的维数 s 与 \(E^{-}\) 的维数 t 对所有满足 a) 中所述条件的直和分解都是相同的；整数 s（相应地，t）是使 \(\Phi\) 在其上的限制为正（相应地，负）且非退化的 E 的子空间 F 的维数的最大值。
\item
  \(\Phi\) 的秩是 \(s + t\) 。
\item
  若 \(\Phi\) 非退化，其指数等于 inf(s, t)（§ 4 第 2 小节定义 2）。
\end{enumerate}

实际上，设 \((x_{i})\) ( \(i = 1, \ldots, n\) ) 是 E 的正交基（§ 6 第 1 小节定理 1）；回顾 \(\Phi(x, x) \in K\) （对一切 \(x \in E\) ）。可以假设 \(\Phi(x_{i}, x_{i}) > 0\) （对 \(i = 1, \ldots, s\) ），\(\Phi(x_{i}, x_{i}) < 0\) （对 \(i = s + 1, \ldots, s + t\) ），\(\Phi(x_{i}, x_{i}) = 0\) （对 \(i = s + t + 1, \ldots, n\) ）。这就证明了 \(a\) )，因为取 \(E^{+}\) （相应地 \(E^{-}\) ）为 \(x_{1}, \ldots, x_{s}\) （相应地 \(x_{s+1}, \ldots, x_{s+t}\) ）生成的子空间，而 \(E^{0}\) 则由 \(x_{s+t+1}, \ldots, x_{n}\) 生成。由此推出 \(b\) )，注意到 \(\Phi(x_{i}, x_{i})\) 具有 \(\rho\bar{\rho}\) 的形式（相应地，

\(-\rho\bar{\rho}) \quad (\rho\in\Lambda^{*}) \quad \text{对} \quad i=1,\ldots,s \quad (\text{resp. } i=s+1,\ldots,s+t).\) 为证明 c)，考察 E 的一个子空间 P，使得 \(\Phi\) 在 P 上的限制是正的且非退化；这时有 \(\mathrm{P}\cap(\mathrm{E}^{-}+\mathrm{E}^{0})=\{0\}\) ，因而和 \(P+E^{-}+E^{0}\) 是直和；我们由此得出 dim \(P\leqslant\dim E^{+}=s\) ，这就证明了 c)。断言 d) 由 a) 立即得出。

最后设 \(\Phi\) 非退化，令 \(q = \inf (s, t)\) ，我们来证明 \(q\) 是 \(\Phi\) 的指数。用 \(b\) ) 的记号，向量 \(e_i + e_{s+i}\) （相应地 \(e_i - e_{s+i}\) ）( \(i = 1, \ldots, q\) ) 生成一个全迷向子空间 F（相应地 F′）。由于 K 的特征为 0，F + F′ 由 \(e_1, \ldots, e_q\) 与 \(e_{s+1}, \ldots, e_{s+q}\) 生成，因而是非迷向的。\(\Phi\) 在子空间 H = (F + F′)⁰ 上的限制于是是正（或负）的且非退化，故 H 不含 \(\neq 0\) 的迷向向量。由于 F⁰ 包含 F 与 H，且 dim(F + H) = codim F′ = codim F，有 F⁰ = F + H。于是与 F 正交的迷向向量 z 必属于 F，因为在直和 F + H 中，z 在 H 中的分量是迷向的。因此 F 是一个极大全迷向子空间，这就证明了 e)。

定义 2. --- 用定理 1 的记号，称偶 \((s,t)\) 为 \(\Phi\) 的符号差。

